% ============================================================
%  Chapter 5 -- Synthesis and Infinite Games
% ============================================================

\chapter{Synthesis and Infinite Games}
\label{ch:synthesis-games}

\aside{Verification asks whether a given
system is correct.  Synthesis asks for a
correct system to be constructed
automatically.}

So far, the system has been part of the input.  We model its behaviours,
compare them with a specification, and look for counterexamples.  The
synthesis problem is more ambitious: given a specification, construct a
finite-state system that satisfies it against every possible behaviour of
the environment.

This chapter presents the game-theoretic view behind Church's synthesis
problem and the B\"uchi--Landweber solution pipeline.


% ============================================================
\section{From Verification to Synthesis}
\label{sec:verification-to-synthesis}
% ============================================================

\begin{definition}[Reactive synthesis]
\label{def:reactive-synthesis}
\textbf{Reactive synthesis} asks whether there exists a finite-state
controller that, for every infinite input sequence produced by the
environment, generates an output sequence satisfying a given
specification.
\end{definition}

The important phrase is ``for every input sequence''.  The controller is
not allowed to choose a friendly environment.  It must react correctly
to all behaviours that the environment may produce.

\begin{definition}[Mealy transducer]
\label{def:mealy-transducer}
A \textbf{Mealy transducer} is a finite-state machine whose output at
each step depends on the current state and the current input symbol.
\end{definition}

This captures causality.  At time $n$, the system may use the input
history up to time $n$, but not future inputs.  A specification that can
only be satisfied by looking into the future is not realizable by a
reactive device.

\begin{example}[Satisfiability is not realizability]
\label{ex:sat-vs-realizable}
The requirement ``output exactly the next input bit'' is satisfiable as a
relation between two infinite sequences: for every input sequence there
exists an output sequence shifted one step earlier.  But it is not
realizable by a causal transducer, because the system would need to know
tomorrow's input today.
\end{example}

\begin{definition}[Church's synthesis problem]
\label{def:church-synthesis}
Given a logical specification $\Phi(X,Y)$, where $X$ are input streams
controlled by the environment and $Y$ are output streams controlled by
the system, decide whether there is a finite-state transducer $M$ such
that for every input stream $x$, the output stream $M(x)$ satisfies
\[
  \Phi(x,M(x)).
\]
If such an $M$ exists, construct it.
\end{definition}


% ============================================================
\section{The Game View}
\label{sec:game-view}
% ============================================================

Synthesis is naturally a game.  The environment chooses inputs, the
system responds with outputs, and the infinite play must satisfy the
specification.

\begin{definition}[Game arena]
\label{def:game-arena}
A \textbf{game arena} is a finite directed graph
\[
  G=(V_E,V_S,E)
\]
whose vertices are split between environment vertices $V_E$ and system
vertices $V_S$.  At an environment vertex, the environment chooses the
next move.  At a system vertex, the system chooses the next move.
\end{definition}

A play is an infinite path through the arena.  The winning condition
selects which infinite paths are winning for the system.

\begin{definition}[Strategy]
\label{def:strategy}
A \textbf{strategy} for a player tells the player which move to take
after each finite history ending in one of that player's vertices.
\end{definition}

Strategies may use memory.  A \textbf{positional} strategy depends only
on the current vertex.  A \textbf{finite-state} strategy depends on the
current vertex plus a finite memory state.

\begin{example}[Request-grant game]
\label{ex:request-grant-game}
The environment may raise a request.  The system may grant.  The
specification is:
\[
  \text{every request is eventually granted.}
\]
This is not a finite-prefix property.  A finite play cannot prove that a
request will never be granted.  The game must be evaluated over infinite
plays.
\end{example}


% ============================================================
\section{Automata-Theoretic Synthesis Pipeline}
\label{sec:synthesis-pipeline}
% ============================================================

The B\"uchi--Landweber solution connects logic, automata, and games.

\begin{formaldetails}[Pipeline]
\begin{enumerate}[label=\textbf{Step \arabic*.}, leftmargin=2cm]
  \item Write the specification in a logic over infinite words, such as
        S1S or a temporal logic fragment.
  \item Translate the specification into an equivalent deterministic
        automaton with an acceptance condition suitable for games, such
        as a M\"uller or parity condition.
  \item Split each round into an environment move followed by a system
        move, obtaining a finite game arena.
  \item Solve the infinite game.
  \item If the system player has a winning strategy, extract a finite
        transducer from that strategy.
\end{enumerate}
\end{formaldetails}

The pipeline differs from model checking in one crucial way.  Model
checking takes a system automaton as input.  Synthesis constructs the
system from a winning strategy.

\begin{definition}[M\"uller condition]
\label{def:muller-condition}
A \textbf{M\"uller condition} specifies a family
$\mathcal{F}\subseteq 2^Q$ of accepting sets.  A run is accepting if
\[
  \inf(\rho)\in\mathcal{F}.
\]
\end{definition}

B\"uchi acceptance asks whether $\inf(\rho)$ intersects a fixed set
$F$.  M\"uller acceptance can inspect the whole set of states seen
infinitely often.  This additional expressiveness is useful when moving
from non-deterministic automata to deterministic automata for synthesis.


% ============================================================
\section{Reachability Games and Attractors}
\label{sec:reachability-games}
% ============================================================

The simplest game-solving pattern is reachability.

\begin{definition}[Reachability game]
\label{def:reachability-game}
In a reachability game, the system wins a play if the play eventually
visits a target set $T\subseteq V$.
\end{definition}

The winning region is computed by an attractor construction.

\begin{definition}[Attractor]
\label{def:attractor}
The system attractor to $T$ is the least set $\mathsf{Attr}(T)$ such
that:
\begin{itemize}[nosep]
  \item $T\subseteq \mathsf{Attr}(T)$;
  \item if $v\in V_S$ has some successor in $\mathsf{Attr}(T)$, then
        $v\in\mathsf{Attr}(T)$;
  \item if $v\in V_E$ has all successors in $\mathsf{Attr}(T)$, then
        $v\in\mathsf{Attr}(T)$.
\end{itemize}
\end{definition}

The two predecessor clauses reflect control.  At a system vertex, one
good successor is enough because the system chooses.  At an environment
vertex, all successors must be good because the environment chooses.

\begin{figure}[ht]
\centering
\begin{tikzpicture}[
  env/.style={circle, draw, thick, fill=orange!12, minimum size=9mm},
  sys/.style={rectangle, draw, thick, rounded corners=2pt, fill=teal!12, minimum size=9mm},
  target/.style={draw, double, thick, fill=accent!16},
  arr/.style={-Stealth, thick},
]
  \node[sys] (s0) {$s_0$};
  \node[env, right=2.2cm of s0] (e1) {$e_1$};
  \node[sys, target, right=2.2cm of e1] (t) {$T$};
  \node[env, below=1.4cm of e1] (bad) {$e_2$};
  \draw[arr] (s0) -- (e1);
  \draw[arr] (e1) -- (t);
  \draw[arr] (e1) -- (bad);
  \draw[arr] (bad) edge[loop right] node {\small stay} (bad);
\end{tikzpicture}
\caption{At environment vertices, every possible successor must be safe for attractor membership.}
\label{fig:attractor-control}
\end{figure}


% ============================================================
\section{M\"uller and Parity Games}
\label{sec:muller-parity-games}
% ============================================================

Reachability is not enough for full synthesis.  Specifications often
depend on what happens infinitely often.

\begin{definition}[Parity condition]
\label{def:parity-condition}
In a parity game, each vertex has a priority in $\mathbb{N}$.  The
winner is determined by the highest priority seen infinitely often:
one player wins when that priority is even, the other when it is odd.
\end{definition}

Parity games are central because many acceptance conditions can be
reduced to them, and they admit memoryless winning strategies.  M\"uller
games can be reduced to parity games by recording enough finite memory
about the latest appearance of states.

\begin{definition}[Latest appearance record]
\label{def:lar}
The \textbf{latest appearance record} (LAR) is a finite memory structure
that keeps the states seen so far ordered by recency.  When a state is
visited again, it is moved to the front.
\end{definition}

The LAR lets a parity condition detect which set of states is seen
infinitely often.  It is a finite-memory encoding of a long-run set.

\begin{formaldetails}[Determinacy fact]
For the finite infinite-duration games used in the synthesis pipeline,
from every vertex one of the two players has a winning strategy.  For
reachability games, this follows directly from the attractor
construction: vertices inside the attractor are winning for the system,
and vertices outside are winning for the environment by avoiding the
attractor forever.

For parity and M\"uller games the proof is deeper.  The course uses the
result as the game-theoretic foundation that makes synthesis decidable:
if the system does not have a winning strategy, then the environment has
a counter-strategy.
\end{formaldetails}

Determinacy is what makes synthesis a decision problem.  If the system
does not have a winning strategy, then the environment has a strategy to
force a violation.


% ============================================================
\section{Formal Verification in Practice}
\label{sec:formal-verification-practice}
% ============================================================

The same automata and game ideas appear in industrial verification.  A
typical safety-critical example is railway interlocking: the controller
must prevent incompatible train movements and unsafe signal settings.

\begin{intermezzo}{Railway Interlocking}
Legacy railway systems were often implemented with relay circuits.  When
such systems are migrated to software, the engineering challenge is not
only to generate code, but to prove that the generated controller
preserves the safety constraints of the original infrastructure.

Model-based design, software model checking, deductive verification, and
reverse engineering of relay logic all fit the same pattern: extract a
finite model, state the property precisely, and use an automated method
to rule out bad behaviours.
\end{intermezzo}

Synthesis pushes this idea further.  Instead of verifying code after it
has been written, we try to generate a correct-by-construction
finite-state controller from the specification.

\begin{summarybox}
\begin{itemize}[nosep]
  \item Synthesis constructs a finite-state system from a specification.
  \item Causality separates realizability from ordinary satisfiability.
  \item Infinite games model the interaction between environment and system.
  \item Attractor computation explains the basic predecessor reasoning.
  \item M\"uller and parity games handle conditions about infinite behaviour.
\end{itemize}
\end{summarybox}

\section*{Exercises}
\addcontentsline{toc}{section}{Exercises}

\begin{chapex}[Causality]
Classify the following specifications as causal or non-causal for a
Mealy transducer: (i) output the current input bit; (ii) output the
previous input bit; (iii) output the next input bit.
\end{chapex}

\begin{chapex}[Attractor iteration]
In \cref{fig:attractor-control}, compute the first attractor iteration
starting from $T$.  Why is $e_1$ not immediately added?
\end{chapex}

\begin{chapex}[Strategy memory]
Give an example of a property where reacting only to the current input
symbol is insufficient, but one bit of memory is enough.
\end{chapex}
